\section{Estimating a modular checksum}
\label{sec:checksum}

We now estimate a sum over a modular subset-sum bin without enumerating
that bin or all Fourier frequencies.  The input residues are arbitrary.
Random phases, rather than a distributional assumption on the input,
will make the moment bounds of Section~\ref{sec:filtering} applicable.

Write $e_P(j)=\exp(2\pi i j/P)$.  For residues
$b_1,\ldots,b_n\in\Z/P\Z$, independent uniform phases
$\phi_1,\ldots,\phi_n\in\Z/P\Z$, and a target residue $\rho$, define
\begin{equation}
 V_\rho=
 \sum_{\substack{I\subseteq[n]\\\sum_{i\in I}b_i=\rho\pmod P}}
 e_P\!\left(\sum_{i\in I}\phi_i\right).
 \label{eq:checksum-definition}
\end{equation}
The phases will be retained so that individual subset contributions can
subsequently be subtracted from this estimate.

\begin{theorem}[Modular checksum estimator]
\label{thm:checksum}
Fix a constant $d>0$.  There is a uniform randomized procedure with the
following guarantee for all sufficiently large $n$.  Its inputs are an
integer $P$ with $|\log P-\eta n|\le2$, arbitrary residues
$b_1,\ldots,b_n$ modulo $P$, and a list $\mathcal T$ of at most $n^d$
target residues.  It draws and retains independent uniform phases
$\phi_1,\ldots,\phi_n$ and either reports failure or returns estimates
$\widehat V_\rho$ for all $\rho\in\mathcal T$.  With probability
$1-o(1)$ it does not fail and
\[
 |\widehat V_\rho-V_\rho|\le\frac18
 \qquad\text{for every }\rho\in\mathcal T.
\]
The $o(1)$ is uniform in $P$, the input residues, and the target list.
On every random execution the procedure takes at most
\begin{equation}
 2^{(\tau+20\eps)n+O(\sqrt n)}\poly(n+|\mathcal T|)
 \label{eq:checksum-time}
\end{equation}
word operations.  No primality assumption on $P$ is needed.
\end{theorem}

The proof has three parts.  We cover every frequency using products of
small random tables, while keeping the two lists of table entries short.
We then count and sample the resulting matches without expanding all
table products.  Finally, inverse-multiplicity weighting and the
filtering inequalities give an exponentially small mean squared error.

\subsection{Frequency patterns and block tables}

Character orthogonality gives the exact identity
\begin{equation}
 V_\rho=\frac1P\sum_{k=0}^{P-1}J_\rho(k),\qquad
 J_\rho(k)=e_P(-k\rho)\prod_{i=1}^n(1+e_P(x_i(k))),\qquad
 x_i(k)=\phi_i+kb_i\pmod P.
 \label{eq:checksum-fourier}
\end{equation}
Indeed, expanding the product leaves one term for every subset, and its
average character is one exactly when its sum is $\rho$ modulo $P$.
In particular,
\begin{equation}
 |J_\rho(k)|^2=\prod_{i=1}^n M(x_i(k))^2,
 \label{eq:checksum-amplitude}
\end{equation}
where residues are identified with their angles in the finite circle.
All functions in this section are the finite-circle functions of
Lemma~\ref{lemma:finite-filtering}, with their finite normalizations.

Partition $[n]$ deterministically into $q=\Theta(\sqrt n)$ blocks
$B_1,\ldots,B_q$ of sizes $d_j=\Theta(\sqrt n)$.  A frequency's
\emph{pattern} records, for each block, rational approximations
$D_j,A_j,E_j$ to the block averages of $D,A,E$ at its coordinates
$x_i(k)$, each with error at most $\eps$.  We choose these approximations
so that
\begin{equation}
 D_j\ge0,\qquad A_j\le D_j,\qquad E_j\le D_j.
 \label{eq:checksum-pattern-order}
\end{equation}
Here and throughout, the approximation rule is fixed and deterministic.
For example, approximate each average to error $\eps/8$, round its
upper endpoint upward for $D_j$, and round its lower endpoint downward
for $A_j,E_j$, on a rational grid of step at most $\eps/4$.
The exact inequalities $D\ge0,D\ge A,D\ge E$ then imply
\eqref{eq:checksum-pattern-order}.  Uniform boundedness of these
functions puts all the rounded values in a fixed finite grid.
Lemma~\ref{lemma:finite-filtering} computes these approximations in
polynomial time per frequency, without summing over $P$ points.
Fix likewise deterministic approximations to the individual scores
$\log f_*,\log g_*$ with error at most $\eps$.

Let $\mathcal P$ be the collection of all blockwise grid triples
satisfying \eqref{eq:checksum-pattern-order}, including patterns that
no frequency attains.  Its size is
\begin{equation}
 Q_n:=|\mathcal P|=2^{O(\sqrt n)}.
 \label{eq:checksum-pattern-count}
\end{equation}
We enumerate this collection, not the $P$ frequency assignments.
The collection and its grid indices are fixed independently of the
phases; only a frequency's membership in a pattern depends on them.
Fix $\pi\in\mathcal P$ and put
\[
 (D_\pi,A_\pi,E_\pi)
   =\frac1n\sum_{j=1}^q d_j(D_j,A_j,E_j),\qquad
 U_\pi=A_\pi+E_\pi-D_\pi,\qquad
 W_\pi=2D_\pi-A_\pi-E_\pi.
\]
For the moment let $0<\alpha\le1$ and $0\le H_\pi\le\eta$
be rational parameters.  The parameter $\alpha$ is the probability of
testing a coordinate, and $H_\pi$ determines the relative lengths of
the two digit lists.  After defining
the tests and their entries, we choose them to keep both lists short.
Choose the integer base
$B_\pi=2^{\lceil(\eta-H_\pi)n\rceil}$ and use the digit ranges
\[
 0\le h<\lceil P/B_\pi\rceil,
 \qquad 0\le\ell<B_\pi.
\]
Their vectors have coordinates
\[
 y_i(h)=\phi_i+hB_\pi b_i\pmod P,
 \qquad z_i(\ell)=-\ell b_i\pmod P.
\]
The pair $(h,\ell)$ represents $k=hB_\pi+\ell$ uniquely, and
$y_i(h)-z_i(\ell)=x_i(k)$.  Pairs representing $k\ge P$ will be
discarded only when contributions are evaluated.  The high and low
digit ranges have sizes $2^{H_\pi n+O(1)}$ and
$2^{(\eta-H_\pi)n+O(1)}$, respectively; together they represent
$2^{\eta n+O(1)}$ pairs.

We test each digit vector separately against random block tables.
This organization is an asymmetric filtering construction; related
filter frameworks also distinguish the two marginal acceptance
probabilities from their joint acceptance probability
\cite{Christiani2017}.  The estimates required here are proved below.
For each block $B_j$, generate
\begin{equation}
 R_j=2^{\lceil(\alpha D_j+4\eps)d_j\rceil}
 \label{eq:checksum-table-count}
\end{equation}
independent tables.  At each coordinate of a table independently, its
label is $*$ with probability $1-\alpha$ and otherwise a uniform
$\theta\in\Z/P\Z$ with probability $\alpha$.  Tables in distinct
blocks and patterns are independent, and all are independent of the
phases.  Distinct table indices are retained even when their labels
coincide.

A label $*$ gives score zero on either side.  A non-$*$ label $\theta$
gives the high and low scores
\[
 \log f_*(y_i(h)-\theta),\qquad
 \log g_*(z_i(\ell)-\theta),
\]
respectively.  Use the fixed approximations to these scores when
testing a vector.  It passes a block table if its average approximate
score is at least $\alpha A_j-3\eps$ on the high side, or at least
$\alpha E_j-3\eps$ on the low side.
A \emph{whole table} is a tuple of one table index from each block.
A vector passes it when it passes all its block tables.  A table
\emph{entry} is a digit index and a whole table it passes.  A
\emph{match} is a high entry and a low entry with the same whole table.
Write $L_\pi$ for the number of matches.  These definitions do not
restrict the represented frequency to have pattern $\pi$.

\subsection{Coverage and list sizes}

The parameters must ensure both that each frequency has a match and
that its two lists of entries can be constructed within the time bound.
The estimates below give the leading exponents for the expected entry
counts as
\[
 H_\pi+\alpha(D_\pi-A_\pi),\qquad
 \eta-H_\pi+\alpha(D_\pi-E_\pi),
\]
apart from $O(\eps)$ rounding slack and subexponential factors.  The
digit ranges contribute $H_\pi$ and $\eta-H_\pi$; the number of tables
and the score thresholds contribute the remaining terms.  Their sum
is $\eta+\alpha W_\pi$, so both exponents can be at most $\tau$
exactly when $\alpha W_\pi\le2\tau-\eta=\Gamma$.
Choose the largest $\alpha\in(0,1]$ satisfying this inequality, and
make the high-side exponent equal to $\tau$:
\begin{equation}
 \alpha_\pi=
 \begin{cases}1,&W_\pi\le\Gamma,\\
 \Gamma/W_\pi,&W_\pi>\Gamma,
 \end{cases}
 \qquad H_\pi=\tau-\alpha_\pi(D_\pi-A_\pi).
 \label{eq:checksum-alpha}
\end{equation}
Use these parameters in the construction, abbreviating
$\alpha=\alpha_\pi$.  Since
$D_\pi-A_\pi,D_\pi-E_\pi\ge0$ and $\alpha W_\pi\le\Gamma$,
\begin{equation}
 \Delta\le H_\pi\le\tau,
 \qquad
 \eta-H_\pi+\alpha(D_\pi-E_\pi)
   =\Delta+\alpha W_\pi\le\tau.
 \label{eq:checksum-digit-balance}
\end{equation}
In particular both digit ranges have size at most
$2^{\tau n+O(1)}$.  We now prove the coverage and size estimates for
this choice.

\begin{lemma}[Coverage and sizes]
\label{lemma:checksum-tables}
For the preceding random construction, the following statements hold.
\begin{enumerate}[label=\textup{(\roman*)}]
\item Conditional on arbitrary phases, the probability that some
$k<P$ has no match in its own pattern is $o(1)$, uniformly in the
phases and residues.
\item In each fixed pattern, the expected number of entries on either
side is at most $2^{(\tau+8\eps)n+O(\sqrt n)}$, and
\begin{equation}
 \E L_\pi\le
 2^{(\eta-\alpha U_\pi+12\eps)n+O(\sqrt n)}.
 \label{eq:checksum-match-expectation}
\end{equation}
The entry estimates hold conditional on the phases.  The match
estimate averages over both phases and tables.
\item With probability $1-o(1)$, all frequencies are covered in their
own patterns and, in every pattern, both entry counts are at most
$2^{(\tau+14\eps)n}$ and
\begin{equation}
 L_\pi\le2^{(\eta-\alpha U_\pi+18\eps)n}.
 \label{eq:checksum-match-good}
\end{equation}
\end{enumerate}
\end{lemma}

\begin{proof}
First fix the phases, $k<P$, its pattern, and one block.  We lower-bound
the chance that one block table accepts both digits representing $k$.
Change the label law as follows: retain the probability of $*$, but
conditional on a non-$*$ label give $w=z_i-\theta$ density $K_{x_i(k)}$
relative to the uniform finite-circle law.  Under this new law the
expected high score, low score, and base-two log likelihood ratio
against the original law are, respectively,
\[
 \alpha A(x_i(k)),\qquad \alpha E(x_i(k)),\qquad
 \alpha D(x_i(k)).
\]
Indeed, a non-$*$ label has likelihood ratio $K_{x_i(k)}(w)$ and
$y_i-\theta=x_i(k)+w$.  The coordinate variables remain independent.
Their absolute values are bounded by a fixed constant, uniformly in
$P,k,\alpha$.  Chebyshev's inequality and a union bound therefore show
that, for sufficiently large block size, all three block averages
are within $\eps$ of their means with probability at least $1/2$
under the new law.

On this event the approximate score averages are at least
$\alpha A_j-3\eps$ and $\alpha E_j-3\eps$: the mean approximation,
the deviation, and the score approximation cost at most $\eps$ each.
The likelihood ratio of the entire block is at most
$2^{(\alpha D_j+2\eps)d_j}$.  Changing back to the original law gives
\begin{equation}
 \Prob(\text{one block table accepts both digits})
 \ge\frac12\,2^{-(\alpha D_j+2\eps)d_j}.
 \label{eq:checksum-coverage-probability}
\end{equation}
Thus the chance that none of its $R_j$ tables accepts both is at most
$\exp(-2^{2\eps d_j-1})$.  Union over the $q$ blocks and the $P$
frequencies is still $o(1)$ because $d_j=\Theta(\sqrt n)$.
If each block has an accepting table, their tuple is a match.  This
proves~\textup{(i)}.

For the size estimates, the number of whole tables is at most
\begin{equation}
 \prod_jR_j\le2^{(\alpha D_\pi+4\eps)n+q}.
 \label{eq:checksum-whole-table-count}
\end{equation}
For any fixed vector and fixed whole-table index tuple, let $S$ be its
total exact score.  Normalization of $f_*$ or $g_*$ and independence
of the labels give $\E 2^S=1$, including the $*$ option.
Passing the tests requires $S\ge(\alpha A_\pi-4\eps)n$ on the high
side, or $S\ge(\alpha E_\pi-4\eps)n$ on the low side.
Markov's inequality, the digit counts, and
\eqref{eq:checksum-whole-table-count} bound the expected entry counts by
\begin{align*}
 2^{(H_\pi+\alpha(D_\pi-A_\pi)+8\eps)n+O(\sqrt n)}
   &=2^{(\tau+8\eps)n+O(\sqrt n)},\\
 2^{(\eta-H_\pi+\alpha(D_\pi-E_\pi)+8\eps)n+O(\sqrt n)}
   &\le2^{(\tau+8\eps)n+O(\sqrt n)}.
\end{align*}
No independence between distinct whole tables is used.

For a fixed digit pair and table-index tuple, let $S_h,S_\ell$ be the
two total exact scores.  This time
\begin{equation}
 \E_{\phi,\text{tables}}2^{S_h+S_\ell}=1.
 \label{eq:checksum-joint-moment}
\end{equation}
To see this, first condition on all labels.  At every non-$*$
coordinate the independent uniform phase makes $y_i-\theta$ uniform,
so averaging its $f_*$ factor gives one.  Averaging the remaining
$g_*$ factors over the table shifts also gives one.
A match requires
$S_h+S_\ell\ge(\alpha(A_\pi+E_\pi)-8\eps)n$.
Markov's inequality, the $2^{\eta n+O(1)}$ digit pairs, and
\eqref{eq:checksum-whole-table-count} prove
\eqref{eq:checksum-match-expectation}.
It is essential here that all digit pairs are counted, before checking
their patterns; the phases in \eqref{eq:checksum-joint-moment} have
their original independent uniform law.

Finally, each bound in~\textup{(iii)} has a factor
$2^{6\eps n-O(\sqrt n)}$ of slack over its expectation.
Markov's inequality and a union bound over
$Q_n=2^{O(\sqrt n)}$ patterns give simultaneous size bounds with
probability $1-o(1)$.  Combine these with~\textup{(i)}.
\end{proof}

\subsection{Constructing and sampling the matches}

The preceding bounds would not be useful if finding the table entries
required visiting every whole table.  We instead construct the lists
block by block.  For each digit vector, test it against the $R_j$ tables
of every block and record its passing index sets.  If any set is empty,
the vector produces no entries.  Otherwise, emit the Cartesian product
of these sets, with the digit index attached to each tuple.
Since all $D_j$ are uniformly bounded,
\[
 \sum_jR_j=2^{O(\sqrt n)}\poly(n).
\]
Consequently the work per digit, apart from emitted entries, is
$2^{O(\sqrt n)}\poly(n)$; a nested counter emits each tuple with
polynomial overhead.  Coincident vectors still retain their different
digit indices.

Sort both entry lists lexicographically by table-index tuple.  If the
high and low group lengths for a tuple are $a_\nu,b_\nu$, then
\begin{equation}
 L_\pi=\sum_\nu a_\nu b_\nu.
 \label{eq:checksum-match-count}
\end{equation}
Prefix sums of these products allow a uniformly random match to be
chosen without enumerating them: choose a uniform integer below
$L_\pi$, locate its group by binary search, and use quotient and
remainder to select one entry on each side.  A zero product contributes
no interval.  The same groups also allow sequential match enumeration.

For any represented frequency $k=hB_\pi+\ell$, define
$\mu_\pi(k)$ to be its number of matches in this pattern procedure.
Its digits are unique, so
\begin{equation}
 \mu_\pi(k)=\prod_{j=1}^q
 \bigl|\{r\in[R_j]:\text{block table }r
       \text{ accepts both }y(h)\text{ and }z(\ell)\}\bigr|.
 \label{eq:checksum-multiplicity}
\end{equation}
We compute this integer directly by retesting the stored block tables,
in $2^{O(\sqrt n)}\poly(n)$ time.  Repeated table labels are counted
with their distinct indices, exactly as in the lists.
The very same deterministic score approximations are used in list
construction and multiplicity computation.  Thus the identity is
exact, not approximate.  Whenever a match is sampled or enumerated,
its multiplicity is positive, even on outcomes where coverage fails.

These operations never scan the frequency range or the set of whole
tables.  They also permit an immediate cap: abort before either entry
list exceeds
\begin{equation}
 C_n=2^{\lceil(\tau+20\eps)n\rceil}
 \label{eq:checksum-cap}
\end{equation}
entries.  Once the lists exist, $L_\pi\le C_n^2$ on every execution,
so the exact match count and its prefix sums have $O(n)$ bits.
Large match counts therefore do not force large enumerations.

\subsection{Direct summation and importance sampling}

Let $\mathcal G$ denote the event in
Lemma~\ref{lemma:checksum-tables}\textup{(iii)}.  We first analyze
exact arithmetic and ideal discrete uniform draws.  Entry and
enumeration caps will be imposed after this analysis; on $\mathcal G$
they do not interrupt it.

For each pattern define its exact contribution
\[
 v_{\pi,\rho}=\frac1P
   \sum_{\substack{k<P\\\text{pattern}(k)=\pi}}J_\rho(k).
\]
If $\alpha U_\pi\ge\Delta$, enumerate every match and give it
contribution zero unless $k<P$ and its deterministic pattern is
$\pi$.  For a retained match use
\begin{equation}
 \frac{J_\rho(k)}{P\mu_\pi(k)}.
 \label{eq:checksum-direct-weight}
\end{equation}
On $\mathcal G$, every frequency of this pattern is present, and its
$\mu_\pi(k)$ contributions sum to $J_\rho(k)/P$.
Also \eqref{eq:checksum-match-good} gives
$L_\pi\le2^{(\tau+18\eps)n}$ in this case.  We cap this enumeration
at $C_n$, aborting before visiting a further match.

If $\alpha U_\pi<\Delta$, take
$N=2^{\lceil\tau n\rceil}$ independent uniform matches with replacement.
When $L_\pi=0$, return zero for this pattern.  Otherwise, average the
sample values
\begin{equation}
 X_{\pi,\rho}=
 \begin{cases}
 \displaystyle\frac{L_\pi}{P\mu_\pi(k)}J_\rho(k),
      &k<P\text{ and }\text{pattern}(k)=\pi,\\[2mm]
 0,&\text{otherwise}.
 \end{cases}
 \label{eq:checksum-sample-weight}
\end{equation}
The inverse-multiplicity correction is the usual principle of counting
each object once across overlapping representations
\cite[Section~5]{KLM1989}.  Here the summands are complex and
we need an absolute error bound, so we give the relevant calculation.
Conditional on phases and tables in $\mathcal G$, a frequency of this
pattern is sampled with probability $\mu_\pi(k)/L_\pi$.  Hence the
sample mean is unbiased for $v_{\pi,\rho}$ and its mean squared error
is at most
\begin{equation}
 \frac1N\,\frac{L_\pi}{P^2}
 \sum_{\substack{k<P\\\text{pattern}(k)=\pi}}|J_\rho(k)|^2.
 \label{eq:checksum-conditional-variance}
\end{equation}
Indeed, its second moment before averaging is the same sum with an
additional factor $1/\mu_\pi(k)\le1$, and independent complex sample
errors have zero cross terms.  If $L_\pi=0$, coverage implies that
there are no frequencies of this pattern, so the zero estimate is exact.

It remains to show that \eqref{eq:checksum-conditional-variance} is
small.  The phases have been conditioned on in that formula; they must
not now be treated as independent under the conditional law of
$\mathcal G$.  Instead, we bound the formula times
$\mathbf1_{\mathcal G}$ pointwise and only then take expectation under
the original phase law.

For a frequency of pattern $\pi$, let $\bar u,\bar v$ be the coordinate
averages of $u,v$ at $x_i(k)$.  The pattern errors imply
\begin{equation}
 |U_\pi-\bar u|\le3\eps,
 \qquad |W_\pi-\bar v|\le4\eps.
 \label{eq:checksum-profile-errors}
\end{equation}
There are two cases within the sampling branch.  When
$W_\pi\le\Gamma$, $\alpha=1$ and $U_\pi<\Delta$.  Thus
\begin{equation}
 -\alpha U_\pi
 \le-\Delta+8.5(\Delta-\bar u)+30\eps.
 \label{eq:checksum-domination-one}
\end{equation}
For example, writing $\delta=\Delta-U_\pi>0$, the right side minus
the left side is at least $7.5\delta+(30-25.5)\eps$.
When $W_\pi>\Gamma$, we have $\alpha aW_\pi=\Delta$, and the sampling
condition implies $aW_\pi-U_\pi>0$.  Since $\alpha\le1<8.3$,
\begin{align}
 -\alpha U_\pi
 &=-\Delta+\alpha(aW_\pi-U_\pi)\notag\\
 &\le-\Delta+8.3(a\bar v-\bar u)+30\eps.
 \label{eq:checksum-domination-two}
\end{align}
Here the rounding loss is at most $8.3(4a+3)\eps<30\eps$.

Set
\[
 R_1(x)=8.5(\Delta-u(x)),\qquad
 R_2(x)=8.3(av(x)-u(x)).
\]
For a fixed sampled pattern, let $r=1$ or $2$ according to the preceding
case.  Combining \eqref{eq:checksum-match-good},
\eqref{eq:checksum-amplitude}, and
\eqref{eq:checksum-domination-one} or
\eqref{eq:checksum-domination-two} gives, for every $k<P$,
\begin{equation}
 \mathbf1_{\mathcal G}\mathbf1_{\{\text{pattern}(k)=\pi\}}
 L_\pi|J_\rho(k)|^2
 \le2^{(\eta-\Delta+48\eps)n}
       \prod_{i=1}^n\bigl(M(x_i(k))^2\,2^{R_r(x_i(k))}\bigr).
 \label{eq:checksum-unconditional-domination}
\end{equation}
The nonnegative right side no longer restricts the pattern or the good
event.  At each fixed $k$, under the original law the $x_i(k)$ are
independent uniform residues, regardless of the values of $b_i$.
Lemma~\ref{lemma:finite-filtering} therefore bounds the expected
product by $2^{0.5169n}$.

Let $E_{\pi,\rho}$ be the error of this pattern's ideal estimate.
Taking conditional expectation over its samples first in
\eqref{eq:checksum-conditional-variance}, and then using
\eqref{eq:checksum-unconditional-domination}, yields
\begin{align}
 \E\bigl[\mathbf1_{\mathcal G}|E_{\pi,\rho}|^2\bigr]
 &\le\frac{2^{(\eta-\Delta+48\eps+0.5169)n}}{NP}\notag\\
 &=2^{(-\tau-\Delta+0.5169+48\eps)n+O(1)}
 \le2^{-0.0007n}
 \label{eq:checksum-error-exponent}
\end{align}
for sufficiently large $n$.  The exponent before the $O(1)$ term is
$-0.00081952$.  Directly summed patterns have zero error on
$\mathcal G$.  Cauchy--Schwarz gives
\[
 \E\!\left[\mathbf1_{\mathcal G}
       \left|\sum_{\pi\in\mathcal P}E_{\pi,\rho}\right|^2\right]
 \le Q_n^2\,2^{-0.0007n}.
\]
Markov's inequality and a union bound over $\mathcal T$ now show that,
with probability $1-o(1)$, $\mathcal G$ holds and every total ideal
estimate has absolute error at most $1/16$.  The subexponential number
of patterns and the polynomial number of targets are both absorbed by
the negative linear exponent.  This conclusion uses no independence
between pattern errors or between target estimates.

\subsection{Resource and precision bounds}

We finish the proof of Theorem~\ref{thm:checksum} by making the
preceding procedure bounded on all histories.  Generate block tables
and digit entries as described, aborting at the entry cap
\eqref{eq:checksum-cap}.  For a direct pattern, abort at the same cap
on enumerated matches.  Sampled patterns use exactly $N$ match draws,
or return zero when there is no match.  The event $\mathcal G$ implies
that none of these caps is reached.  On any other history, the caps
still bound the work.

Each grid value has a fixed rational denominator.  Thus the weighted
averages and $\alpha$ have numerators and denominators of $O(\log n)$
bits, even for unattained patterns.  The ceilings in the digit and
table counts are computed exactly by rational arithmetic.
For each pattern, digit processing costs
$2^{\tau n+O(\sqrt n)}\poly(n)$ before the emitted entries.
Sorting, prefix sums, and random indexing have polynomial cost per
entry or sample.  Each valid contribution requires at most
$2^{O(\sqrt n)}\poly(n)$ work for its multiplicity, and a polynomial
factor for the target list.  All tuple keys, match counts, and
multiplicities have $O(n)$ bits: for multiplicities use
\eqref{eq:checksum-whole-table-count}, and for match counts on retained
lists use $L_\pi\le C_n^2$.
Multiplication by the $Q_n$ patterns gives
\eqref{eq:checksum-time}.

The coarse fixed approximations in pattern and acceptance tests define
the lists themselves; they need not be refined later.  The complex
contributions require greater accuracy.  On every nonaborted history
there are only $2^{O(n)}\poly(|\mathcal T|)$ of them, each formed from
at most $n+1$ factors of modulus at most two.  The rational weights in
\eqref{eq:checksum-direct-weight} and
\eqref{eq:checksum-sample-weight} have $O(n)$-bit numerators and
denominators and magnitude at most $2^{O(n)}$.
Section~\ref{sec:finite-arithmetic}, under
\emph{Complex evaluations and summation}, implements precisely such
sums by evaluating each phase to error $2^{-2n^2}$ and rounding each
contribution to $n^2$ fractional bits on a common dyadic grid.  Its
total error here is
$2^{-n^2+O(n)}\poly(|\mathcal T|)=o(1)$, hence less than $1/16$
for sufficiently large $n$, with polynomial work per contribution.

The rational probabilities in the label choices, finite uniform draws,
and phase evaluations admit bounded implementations on the specified
word RAM; Section~\ref{sec:implementation} gives these implementations
and their coupling to the ideal draws.  Their additional abort
probability is $o(1)$, and their cost is polynomial per operation just
counted.  The retained phases and block labels are ordinary writable
data; the procedure requires no access to unstored randomness.
The ideal error bound $1/16$, the arithmetic error $1/16$, and the
vanishing probabilities of cap or sampling failure prove
Theorem~\ref{thm:checksum}.
